\hypertarget{appendix-6854d0e2218e}{}
\section{Floating-point encoding and decoding}\label{appendix-6854d0e2218e-floating-point-encoding-and-decoding}

\textbf{Curriculum depth:} Introduced · \textbf{First seen:} C004 · \textbf{Latest development:} C004

\subsection{Definition and mechanism}\label{appendix-6854d0e2218e-definition-and-mechanism}

Encoding chooses fields for a value; decoding classifies fields and reconstructs their meaning. 10.625 normalizes to 1.010101₂×2³, giving 412A0000. C1540000 decodes to −13.25.

\subsection{Prerequisites}\label{appendix-6854d0e2218e-prerequisites}

\begin{itemize}
\tightlist
\item
  \hyperref[appendix-2fae058c82e1]{Exponent bias} --- Stored exponents must be translated to effective powers before reconstructing value or spacing.
\item
  \hyperref[appendix-a24fd710cf67]{Significand and hidden bit} --- Fraction weights and effective precision determine value and the smallest retained position.
\item
  \hyperref[appendix-6b8075763724]{Hexadecimal} --- Four-bit groups make complete representation words independently inspectable.
\end{itemize}

\subsection{Used by}\label{appendix-6854d0e2218e-used-by}

Connect later dependent concepts here as their lessons arrive.

\subsection{Recall and next depth}\label{appendix-6854d0e2218e-recall-and-next-depth}

Reconstruct the concrete example in the definition and identify its validity assumptions.

\textbf{Pending:} Inexact conversion requires rounding; the exact worked examples do not explain every decimal conversion.

\subsection{Source and navigation}\label{appendix-6854d0e2218e-source-and-navigation}

\hyperref[01-foundations-i-ch001-a-stretching-ruler-for-real-valued-quantities]{M004 teaching chapter} · \hyperref[appendix-fd090d982a15-m004---parts-61-76]{Per-part definitions} · \href{../02-Foundations-II/main.pdf\#nameddest=appendix-4ad47bdd269a}{Dependency Map (02-Foundations-II, appendix-4ad47bdd269a)}
